\documentstyle[%


righttag%


,11pt%


,amssymb%


,russian%               remove in Eng.


,izvru%                 Set izven% in Eng.


]{amsart}





%


%% Adjust \hspace in Eng.


\newcommand{\mes}{\operatorname{mes}}


\newcommand{\myfont}{\series{m}\shape{n}\selectfont}


\newcommand{\tts}[1]{}





\theoremstyle{plain}


\theorembodyfont{\it}


\newtheorem{thmnonum}{\hskip\parindent Теорема}


\renewcommand{\thethmnonum}{}





\theoremstyle{definition}


\theorembodyfont{\rm}


\newtheorem{notenonum}{\hskip\parindent Замечание}


\newtheorem{examnonum}{\hskip\parindent Пример}


\renewcommand{\thenotenonum}{}


\renewcommand{\theexamnonum}{}





%\hoffset=-15mm%                  For Eng.


 


\setcounter{page}{19}


\god{1998}


\nomer{5~(432)} %                        Remove in Eng.


%\nomer{Vol.\,42, No.\,5}%               Set in Eng.


%\str{\pageref{begin}--\pageref{end}}%  Set in Eng.


\udk{517.983.26}


\author{\it Н.А.\,Ерзакова}


% NB:   ^^ remove in Eng.


\title{Константа Макенхаупта и мера некомпактности\\


оператора вложения пространств Соболева}





\date{}


%\pagestyle{myheadings}%         Set in Eng.





\pagestyle{plain} %              Remove in Eng.


\begin{document}





%\thispagestyle{plain}%          Set in Eng.





\maketitle


\label{begin}








\section*{Введение}





Верхняя $\chi$-норма --- это характеристика степени некомпактности


линейного оператора.


Во многих случаях (напр., [1]--[4]), когда она достаточно мала,


результаты обобщаются с компактных операторов на некомпактные.


Однако бывают ситуации, как в случае с оператором вложения


пространств Соболева на области с нерегулярной границей, когда вычисление


$\chi$-нормы довольно сложно.





В данной работе получены оценки для верхней $\chi$-нормы


оператора вложения специального класса пространств Соболева,


включающего пространства Соболева на области с так называемым


``обобщенным хребтом'' [5], [6].





Пусть $Q$ --- область в $R^n$, $W^1_p(Q)=W^1_p$ --- пространство Соболева,


$L_q(Q)=L_q$ --- пространство Лебега, $1<p,q<\infty$, $\mes Q<\infty$,


$L^1_p(Q)=L^1_p$ --- пространство всех обобщенных функций на $Q$


с полунормой $\|f\|_{L^1_p}=\|\nabla f\|_{L_p}$, $C$ --- совокупность


всех постоянных функций на $Q$, $P_D$ --- оператор умножения на


характеристическую


функцию измеримого подмножества $D\subset Q$.





\begin{defn}[{\myfont{}[1], [2], [7]}] Пусть $E$ --- банахово пространство,


а $\Omega$ --- ограниченное


подмножество $E$. Мерой некомпактности Хаусдорфа


$\chi_E(\Omega)=\chi(\Omega)$  множества $\Omega$


называется инфимум всех $\varepsilon>0$, при которых $\Omega$ имеет в $E$


конечную $\varepsilon$-сеть.


\end{defn}





\begin{defn}[{\myfont{}[7]}] Пусть $G$ и $E$ --- банаховы пространства,


$A$ --- линейный ограниченный оператор из $G$


в $E$. Верхней $\chi$-нормой оператора $A$ называется величина


$\|A\|^{(\chi)}=\chi_E(AS_1)$, где $S_1$ --- единичная сфера пространства


$G$.


\end{defn}





Пусть $\mu$ --- неотрицательная борелевская мера на $J=[a,b)$


и $\check\mu$ --- абсолютно непрерывная часть $\mu$, $1<p\le q<\infty$.


Под константой Макенхаупта подразумевается величина


\begin{equation}


\lim_{\delta\to 0} N_{\delta},


\end{equation}


где


$$


N_{\delta}= \sup\limits_{r\in J_{\delta}}\biggl\{[\mu(b-)-\mu(r)]^{1/q}


\biggl(\,\int_c^r(d\check\mu/dt )^{-1/(p-1)}dt\biggr)^{(p-1)/p}\biggr\}


$$


(напр., [8], теорема 1, гл.\,1, \S\,1.3, c.\,42).





Всюду будем предполагать наличие вложения


$L_p^1/C$  в $L_q/C$ для $1<p\le q<\infty$,


под которым понимается


справедливость неравенства Пуанкаре


\begin{equation}


\|f-\overline f\|_{L_q}\le k \|\nabla f\|_{L_p},


\end{equation}


где


$$


\overline f =\frac{1}{\mes Q}\;\int_Qf(x)dx


$$


([8], лемма 2, гл.\,3, \S\,3.2, c.\,143; [6], неравенство


2.18).





Будем предполагать, что область $Q$ удовлетворяет условиям:





1. для некоторого интервала $J=[a,b)$ существует равномерно локально


липшицево отображение $\tau :Q\to J$, т.\,е. для каждого


$x\in Q$ найдется такая окрестность $V(x)$,


что


\begin{equation}


|\tau (x)-\tau (y)|\le \gamma | x-y|


\end{equation}


для всех $y\in V(x)$, причем $Q\backslash Q_{\delta}$


удовлетворяет условию конуса, где $Q_{\delta}=\tau^{-1}(J_{\delta})$,


а $J_{\delta}=[1/\delta ,b)$ при $b=\infty$, $J_{\delta}=[b-\delta ,b)$~


при $b<\infty$;





2. если


\begin{equation}


\mu (t)=\mes\{x\in Q: \tau (x)\in [a,t)\},\ \;


\mu (d-)-\mu (c)=\mu (c,d)\ \; (a\le c<d\le b),


\end{equation}


$L^p(J,d\mu)=L^p$ ---


пространство функций, локально липшицевых на $J$


и таких, что


$$


\|F\|_{L^p}=\biggl(\,\int_J|F|^p d\mu\biggr)^{1/p}<\infty,


$$


$L^{1,p,q}(J,d\mu )=L^{1,p,q}$ --- множество $F\in L^q$,


для которых $F'\in L^p$, то существует отображение


$M:L^1_p\cap C^1(Q)\to L^{1,p,q}$, удовлетворяющее неравенству


\begin{equation}


\|(Mf)'\|_{L^p}\le\rho\|\nabla f\|_{L_p};


\end{equation}





3. для оператора $TF=F(\tau(x))$, действующего из $L^{1,p,q}$


в $L^1_p$, справедливо равенство


\begin{equation}


\lim_{\delta\to 0} \sup\limits_{f\in S\cap C^1(Q)}


\|P_{Q_{\delta}}(f-TMf)\|_{L_q}=0\quad


(S=\{f\in L^1_p:\|\nabla f\|_{L_p}=1\}).


\end{equation}





\begin{notenonum}


Как следует из лемм 4.2--4.5, теорем 4.6--4.7 работы [6],


области с так называемым ``обобщенным хребтом''


удовлетворяют предположениям теоремы


для случая $p=q$ (только этот случай рассмотрен


в [6]).


\end{notenonum}





Напомним ([6], определение 4.1), что область $Q$ называется областью


с ``обобщенным хребтом'' (generalized ridged domain),


если существуют функции $u$, $\rho$, $\tau$, интервал $J=[a,b)$


$(b\le\infty)$ и положительные константы $\alpha$, $\beta$,


$\delta_1$, $\gamma$ такие, что выполнены условия:


\begin{itemize}


\item[1)] $u:J\to Q$, $\rho:J\to R^+=(0,\infty)$ удовлетворяют


условию Липшица;


\item[2)] $\tau:Q\to J$ равномерно локально удовлетворяет условию


Липшица, т.\,е. для всех $x\in Q$ найдется такая окрестность


$V(x)$, что


$|\tau(x)-\tau(y)| \le\gamma \| x- y\|$ для всех $y\in V(x)$;


\item[3)] $\|x-u(\tau(x))\|\le \alpha\rho(\tau(x))$ для всех $x\in Q$;


\item[4)] $r_n\{\|u'(t)\|+\|\rho'(t)\|\}\le\beta$


для всех


$t\in J$, где $r_n$ --- наилучшая константа в неравенстве


$|(a,b)|\le r_n \|a\|\,\|b\|$, а $\|\cdot\|$ ---


норма в $R^n$;


\item[5)] для


$B_t:=B(u(t),\rho(t))$ и $S(x):=\{y:sy+(1-s)x\in Q$


для всех $0\le s\le 1\}$ имеет место включение


$B_{\tau(x)}\subset Q$ и, кроме того, $S(x)\cap B_{\tau(x)}$


содержит такой шар $B(x)$, что


$\mes(B(x))/\mes(B_{\tau(x)})\ge\delta_1>0$.


\end{itemize}





Тогда кривая $t\to u(t):J\to Q$ называется


``обобщенным хребтом''.





Отсюда видно, что предположение (3) совпадает


с предположением 2) области с ``обобщенным хребтом''. Заметим


также, что в силу предположения (6) отображение


$T$ можно интерпретировать как обратное отображение к $M$


на той части, где происходит ухудшение области.





\section{Формулировка основного результата. Комментарии. Пример}





Для верхней $\chi$-нормы оператора вложения


$I:L^1_p/C\to L_q/C$


выделенного класса пространств Соболева


в пространства Лебега установим связь с константой Макенхаупта.





\begin{thmnonum} Пусть выполнены предположения {\em (1)--(6)}. Тогда


справедливо неравенство


$$


1/\gamma \lim_{\delta\to 0} N_{\delta}\le


\|I\|^{(\chi)}\le \rho q^{1/q}(q/(q-1))^{(p-1)/p}


\lim_{\delta\to 0}N_{\delta}.


$$





В частности, оператор вложения $E:W_p^1\to L_q$ компактен, если


и только если


$$


\lim_{\delta\to 0} N_{\delta}=0.


$$


\end{thmnonum}





\begin{notenonum}


Этот результат обобщает теорему 3 [9],


из него вытекает следствие 5.2 [6] (критерий компактности


оператора вложения в случае $p=q$).


Оценки


степени некомпактности для оператора вложения,


полученные в теореме для $1<p\le q<\infty$


(при выполнении (2)), являются важными для приложений


([1]--[4]).


\end{notenonum}





Понятие ``обобщенного хребта'' не используется из-за


сложности применения этого понятия при $p<q$,


а потому вопрос о его существовании не исследуется.


Кроме того, в нижеследующем примере область $Q$ и $\tau$ удовлетворяют


предположениям теоремы, однако условие 5) определения области с


``обобщенным хребтом'' для них не выполнено.





\begin{examnonum}


Рассмотрим область


$$


Q\subset R^2,\ \; Q=\{(x_1,x_2): 0<x_2<\infty,\


0<x_1<e^{-\xi x_2},\ \; \xi =\text{const}>0\}.


$$


Мера Лебега ее конечна, хотя область и


неограничена, а потому пространства $L_q(Q)$ включают в себя


постоянные функции, поскольку $c^q\mes Q<\infty$.


Область $Q$ удовлетворяет и остальным предположениям теоремы, например,


при $p=q=2$.


Действительно, в качестве $J$ рассмотрим интервал $[0,\infty)$.


Положим $\tau :(x_1,x_2)\to x_2$. Очевидно, $\tau$ удовлетворяет


предположению (3), причем $\gamma=1$.


\end{examnonum}





Для $f\in L^1_2\cap C^1(Q)$ пусть


$$


Mf(t)=e^{\xi t}\int_0^{e^{-\xi t}}f(\tau,t)d\tau.


$$


Нетрудно проверить, что для любого $(x_1,x_2)\in Q$


$$


|f(x_1,x_2)-TMf(x_1,x_2)|\le e^{-\xi x_2}\|\nabla f\|_{L_2}/\xi,


$$


откуда следует справедливость предположения (6).


Непосредственно устанавливается справедливость всех остальных предположений


теоремы.





Вычисляем


\begin{gather*}


\mu(t)=\int_0^te^{-\xi \tau}d\tau = (1-e^{-\xi t})/\xi,\\


\lim_{\delta\to 0}N_{\delta} = \lim_{\delta\to 0}


\sup\limits_{r>1/\delta}(e^{-\xi r}/\xi)^{1/2}\biggl(\,\int_{1/\delta}^r


{(e^{-\xi \tau})}^{-1}d\tau\biggr)^{1/2}=1/\xi.


\end{gather*}


Отсюда


$\|I\|^{(\chi)}\le 2/\xi$ и $\|I\|^{(\chi)}>0$, т.\,е.


$I:L^1_2/C\to L_2/C$ не является компактным, в то же время


верхняя $\chi$-норма может быть сколь угодно малой


при $\xi\to\infty$.





Более сложные примеры областей, удовлетворяющих предположениям


теоремы, можно найти в [5], [6].





\section{Доказательство основного результата}





Доказательство теоремы проведем с помощью двух лемм.





\begin{lem} Пусть выполнены {\em (2)--(6)}.


Тогда


$$


\hspace{\parindent}1.


\hspace{5.8cm} \|\nabla(TF)\|_{L_p}\le \gamma \|F'\|_{L^p},


\hspace{5.8cm} (7)


$$





$2.$ $\lim\limits_{\delta\to 0}\mes Q_{\delta}= 0,$








{\em 3.} если


$I^0=\lim\limits_{\delta\to 0} \sup\limits_{F\in JS}


\|P_{J_{\delta}}(F-F(c))\|_{L^q}$


и


$I_0=\lim\limits_{\delta\to 0} \sup\limits_{F\in JS}\|P_{J_{\delta}}


(F-F(c_0))\|_{L^q}$,


где $JS=\{F\in L^{1,p,q}:\|F'\|_{L^p}=1\}$, $c=b-\delta$


при $b<\infty$


и $c=1/\delta$ при $b=\infty$, $c_0\in \tau(Q)$ ---


произвольная, не зависящая от


$\delta$, точка из интервала, то


$\|I\|^{(\chi)}<\rho I^0$


и $I_0\le \|I\|^{(\chi)}/\gamma$


для любого $c_0$, где


$I:L^1_p/C\to L_q/C$ ---


оператор вложения.


\end{lem}





\begin{notenonum} В определении величин $I^0$ и $I_0$


под $F(c)$ и 


$F(c_0)$ понимается значение непрерывной функции в точке из интервала,


поскольку все функции из $L^{1,p,q}$ непрерывны. Точка $c$


зависит от $\delta$ в отличие от точки $c_0$.


Заметим также, что на классе функций $F+C$, где $C$ --- множество


всех постоянных функций на интервале, разность $F-F(c)$


постоянна для любой точки $c$.


\end{notenonum}








{\bf Доказательство.} 1. Неравенство (7) по лемме 4.4 


[6] вытекает из предположения (3).





2. Имеем $\mes(Q\backslash P)<\varepsilon$ для произвольного


$\varepsilon >0$


и некоторого замкнутого множества $P\subset Q$.


Из непрерывности $\tau$ следует замкнутость


$\tau(P)\subset [\alpha,\beta]\subset J$.


Отсюда найдется такое $\delta>0$, что


$P\subset Q\backslash Q_{\delta}$,


т.\,е. справедливо утверждение 2.





3. Покажем неравенство


$\|I\|^{(\chi)}\le \rho I^0$. Из предположения (6)


следует, что для произвольного $\varepsilon>0$ найдется такое


$\delta >0$, что


\addtocounter{equation}{1}


\begin{equation}


\|P_{Q_{\delta}}(f-TMf)\|_{L_q}<\varepsilon


\end{equation}


для всех $f\in S\cap C^1(Q)$. По определению $I^0$


без ограничения общности считаем справедливым неравенство


\begin{equation}


\|P_{J_{\delta}}(F-F(c))\|_{L^q}\le (I^0+\varepsilon)\|F'\|_{L^p},


\end{equation}


где $c=1/\delta$, если


$b=\infty$, и $c=b-\delta$, если $b<\infty$.





Множество


$U:=\{f-\overline f: f\in S\cap C^1(Q)\}$ в силу


неравенства (2) равномерно ограничено в


$L_q\cap L^1_p$. Отсюда, принимая во внимание (5) и (8), имеем


равномерную ограниченность


$\{P_{J_{\delta}}Mf:f\in U\}$


в $L^{1,p,q}$, что влечет по теореме Соболева равномерную ограниченность


в $L^{\infty}$ множества постоянных


$\{Mf(c):f\in U$, где $c=1/\delta$ при $b=\infty$, $c=b-\delta$ при


$b<\infty\}$.





Оценим расстояние между $U$ и вполне ограниченным в $L_q$


множеством


$\{P_{Q\backslash Q_{\delta}}f + Mf(c)P_{Q_{\delta}}1:f\in U\}$.


Из (5), (8) и (9) имеем для всех $f\in U$


\begin{multline*}


\|f-P_{Q\backslash Q_{\delta}}f-Mf(c)P_{Q_{\delta}}1\|_{L_q}\le


\varepsilon + \|P_{Q_{\delta}}(TMf-Mf(c))\|_{L_q}= \\


=\varepsilon+\|P_{J_{\delta}}


(Mf-Mf(c))\|_{L^q}\le \rho (I^0+\varepsilon)+\varepsilon.


\end{multline*}


Отсюда, учитывая, что по теореме 1


([8], гл.\,1, \S\,1.1, c.\,16) $L^1_p\cap C^{\infty}(Q)$ плотно в $L^1_p$


и $\|I\|^{(\chi)}=\chi_{L_q/C}(IS)=\chi_{L_q}(IU)$,


получим необходимое неравенство.





Докажем последнее утверждение леммы 1, а именно, неравенство


$I_0\le\|I\|^{(\chi)}/\gamma$ для всех $c_0$.


Очевидно,


\begin{equation}


I_0=\lim_{\delta\to 0} \sup\limits_{F\in JS} \|P_{J_{\delta}}


((F-\overline F)-(F-\overline F)(c_0))\|_{L^q}.


\end{equation}


По формуле (2) из работы [3] имеем


\begin{equation}


\|I\|^{(\chi)} \ge \lim_{\delta\to 0} \sup\limits_{TF\in S}


\|P_{Q_{\delta}}(TF-\overline{TF})\|_{L_q}\ge


(1/\gamma)\lim_{\delta\to 0} \sup\limits_{TF\in S} \|P_{J_{\delta}}


(F-\overline{F})\|_{L^q}.


\end{equation}


Из определения $T$, неравенств (7) и (9) вытекает


$$


\|F-\overline F\|_{L^q}=\|TF-\overline{TF}\|_{L_q}\le k\|\nabla TF\|_{L_p}


\le k\gamma\|F'\|_{L^p},


$$


откуда следует равномерная ограниченность


$\{F-\overline F: F\in JS\}$


в $L^{1,p,q}$.


Применяя теорему Соболева, получим равномерную ограниченность постоянных


$\{(F-\overline F)(c_0): F\in JS\}$ в $L^{\infty}$


и, следовательно,


$$


\lim_{\delta\to 0}\sup\limits_{F\in JS} \|P_{J_{\delta}}


(F-\overline F)(c_0)\|_{L^q}= 0,


$$


что в силу (10) и (11) завершает доказательство.








\begin{lem} Пусть $\mu$ --- неотрицательная борелевская мера на $J=[a,b)$


и $\check\mu$ --- абсолютно непрерывная часть $\mu$, $1<p\le q<\infty$,


$N_\delta$ определено в $(1)$.





Тогда


$$


I^0\le q^{1/q}(q/(q-1))^{(p-1)/p}\lim_{\delta\to 0}N_{\delta}


$$


и


$$


\lim_{\delta\to 0} N_{\delta}\le I_0.


$$


\end{lem}





{\bf Доказательство} дословно повторяет рассуждения, проводимые при 


доказательстве неравенства Харди ([8], теорема 1.3.1).





Применяя лемму 2 к мере $\mu$, порожденной неубывающей функцией $\mu(t)$


из (4), являющейся в силу построения неотрицательной борелевской


мерой на $J=[a,b)$, из леммы 1 получим утверждение


теоремы.





\begin{thebibliography}{9}





\bibitem{1}


Садовский Б.Н. {\em Предельно компактные и уплотняющие операторы} //


УМН. -- 1972. -- Т.\,27. -- \No\,1. -- С.\,81--146.


\bibitem{2}


Ахмеров Р.Р., Каменский М.И., Потапов А.С. и др. {\em Меры некомпактности и 


уплотняющие операторы}. -- Новосибирск: Наука, 1986. -- 264~с.


\bibitem{3}


Ерзакова Н.А. {\em Некоторые приложения мер некомпактности к уравнениям


математической физики} // Сб. научн. тр. НИИ КТ. --


Хабаровск, 1994. -- Вып.\,1. -- C.\,38--49.


\bibitem{4}


Yerzakova N.A. {\em The measure of noncompactness of Sobolev


embeddings}// Integr. Equat. and Oper. Theory. -- 1994. --


V.\,19. -- \No\,3. -- P.\,349--359.


\bibitem{5}


Amick C.J. {\em Some remarks on Rellich's theorem and the Poincar\'e


inequality} // J. London Math. Soc. -- 1978. -- V.\,18.


-- \No\,1. -- P.\,81--93.


\bibitem{6}


Evans W.D., Harris D.J. {\em Sobolev embedding for generalized ridged


domains} // Proc. London Math. Soc. -- 1987. -- V.\,54.


-- \No\,1. -- P.\,141--175.


\bibitem{7}


Гольденштейн Л.С., Гохберг И.Ц., Маркус А.С. {\em Исследование


некоторых свойств линейных ограниченных операторов в связи с их


$q$-нормой} // Учен. зап. Кишиневск. ун-та. Сер. физ.-матем. наук. --


1957. -- Т.\,29. -- С.\,29--36.


\bibitem{8}


Мазья В.Г. {\em Пространства С.Л.\,Соболева}. -- Л.: Изд-во ЛГУ,


1985. -- 415 c.


\bibitem{9}


Ерзакова Н.А. {\em Критерии компактности оператора вложения


пространств Соболева} // Сб. научн. тр. НИИ КТ. -- Хабаровск,


1993. -- Вып.\,1. -- С.\,103--109.





\end{thebibliography}





\vskip 2cm





\postup{\it Хабаровский государственный}{$\qquad$\it Поступили}


\postup{\it технический университет}{{\it первый вариант} 14.11.1995}


\postup{}{{\it окончательный вариант} 12.11.1996}





\label{end}


\end{document}


